\section{Agent Workflow 深度剖析}

\subsection{发现与理解}

每个 agent 都以 intent detection start：把 natural language 输入映射到 structured goal。团队会为常见 intent 预定义 slot template。Burak 介绍的实践是先用 small model 做 classification，再用 large model 做 planning。

\subsection{规划与工具选择}

planner 既可以是 chain-of-thought prompt，也可以是 learned policy。Burak 展示了一个 hybrid approach：先用 policy 生成 plan tree，再用 LLM 依次展开。关键在于 tool selection module：它要考虑 tool current status（是否健康）、estimated latency 以及 user context 。

\subsection{执行与复核}

Executor 其实就是 tool orchestration layer。除了调用 API，还要 handle retries、fallbacks 和 confirm-with-user。Burak 强调：对于金融类操作，一定要加入 explicit confirmation step，并把决策记录在 audit log 中。

\subsection{收尾与迭代}

Workflow 的最后一个阶段是总结和 logging：生成 summary 给用户、更新 memory、把 all interactions 记录到 observability pipeline，并根据 failure tag 触发 training pipeline。

\subsection{本章小结}

把 agent workflow 拆解成发现、规划、执行、收尾四步，有助于在每个阶段落地指标、监控与审查。

